
%============================================================
% MATH163 — Comprehensive but Concise Revision Navigator (v3)
%============================================================
\documentclass[11pt]{article}

% Page + Typography
\usepackage[top=2.2cm,bottom=2.2cm,left=2.2cm,right=2.2cm]{geometry}
\usepackage{microtype}
\usepackage{mathpazo}
\linespread{1.06}
\setlength{\parskip}{0.55em}
\setlength{\parindent}{0pt}

% Maths & Tables
\usepackage{amsmath,amssymb,mathtools}
\usepackage{booktabs,tabularx,array}
\usepackage{enumitem}
\usepackage{pifont}
\setlist{nosep,leftmargin=*}

% Colors & Boxes
\usepackage[dvipsnames]{xcolor}
\definecolor{Navy}{HTML}{1B3A57}
\definecolor{Accent}{HTML}{0B6FA4}
\definecolor{Soft}{HTML}{F4F7FB}
\usepackage[skins,most]{tcolorbox}
\tcbset{enhanced,boxsep=2pt,arc=1.5pt,outer arc=1.5pt}

% Titles
\usepackage{titlesec}
\titleformat{\section}{\Large\bfseries\color{Navy}}{}{0em}{}
\titleformat{\subsection}{\large\bfseries\color{Navy}}{}{0em}{}

% Hyperref
\usepackage[hidelinks]{hyperref}

% Boxes
\newtcolorbox{keybox}{colback=Soft,colframe=Accent,left=6pt,right=6pt,top=8pt,bottom=8pt}
\newtcolorbox{notebox}{colback=white,colframe=Accent,left=6pt,right=6pt,top=6pt,bottom=6pt}

\usepackage{rotating}
\usepackage{array}
\setlength{\extrarowheight}{6pt} % extra vertical spacing in rows

%============================================================
\begin{document}

{\LARGE\bfseries MATH163 — Concise Revision Navigator}\par
{\large \textcolor{Navy}{A complete but compact map of ideas for fast revision.}}
\vspace{0.7em}
\hrule
\vspace{1em}

%============================================================
\section{Foundations: Random Variables, PMF/PDF, CDF}

\begin{keybox}
Random variables model \textbf{distributions of possible values}. A \textbf{parameter} describes a population; a \textbf{statistic} summarises a sample; an \textbf{estimator} is a sample-based function estimating a parameter.
\end{keybox}

\subsection*{PMF, PDF, CDF}
For a discrete $X$ with PMF $p_X(x)$:
\[
\Pr(X=x)=p_X(x), \qquad \sum_x p_X(x)=1.
\]
For a continuous $Y$ with PDF $f_Y(y)$:
\[
\Pr(a<Y<b)=\int_a^b f_Y(t)\,dt, \qquad \int_{-\infty}^{\infty} f_Y(y)\,dy=1.
\]
CDFs:
\[
F_X(x)=\Pr(X\le x)=\sum_{i\le x} p_X(i),
\qquad
F_Y(y)=\Pr(Y\le y)=\int_{-\infty}^y f_Y(t)\,dt.
\]
CDFs are monotone, defined on all real numbers; continuous RVs have continuous CDFs.

\subsection*{Expectations and variance}
\[
\mathbb E[X]=\sum_x x p_X(x),
\quad
\mathrm{Var}(X)=\sum_x (x-\mathbb E[X])^2p_X(x)=\mathbb E[X^2]-\left(\mathbb E[X]\right)^2.
\]
\[
\mathbb E[Y]=\int y f_Y(y)dy,
\quad
\mathrm{Var}(Y)=\int (y-\mathbb E[Y])^2 f_Y(y)dy=\mathbb E[Y^2]-\left(\mathbb E[Y]\right)^2.
\]
Median $m$ satisfies $F(m)=1/2$.

%============================================================
\section{Transformations, Linearity, Covariance}

\subsection*{Linearity}
\[
\mathbb E[aX+b]=a\mathbb E[X]+b,
\qquad
\mathrm{Var}(aX+b)=a^2\mathrm{Var}(X).
\]
If $X,Y$ independent:
\[
\mathbb E[XY]=\mathbb E[X]\mathbb E[Y],
\qquad
\mathrm{Var}(X+Y)=\mathrm{Var}(X)+\mathrm{Var}(Y).
\]
Covariance:
\[
\operatorname{Cov}(X,Y)=\mathbb E[XY]-\mathbb E[X]\mathbb E[Y].
\]
General variance identity:
\[
\mathrm{Var}(X+Y)=\mathrm{Var}(X)+\mathrm{Var}(Y)+2\operatorname{Cov}(X,Y).
\]

\textit{Note:} Covariance $=0$ does not imply independence except in special families (e.g.\\ multivariate normal).

\subsection*{Transformations and the 5-step method}
If $Y=G(X)$ with $G$ monotone:
\[
f_Y(y)=f_X(G^{-1}(y))\,\left|\frac{d}{dy} G^{-1}(y)\right|.
\]
\begin{notebox}
\textbf{5‑step method used in lectures:} (1) write CDF of $X$, (2) express $F_Y(y)=\Pr(G(X)\le y)$, (3) invert, (4) apply CDF of $X$, (5) differentiate.
\end{notebox}

%============================================================
%\section{Distribution Reference Sheet}
%Distributions marked with \ding{72} require memorised means and variances.

\begin{sidewaystable}%[Ht]
	\centering
	\small
	\renewcommand{\arraystretch}{1.6}
	\begin{tabular}{l l l l l}
		\toprule
		\textbf{Distribution} & \textbf{Parameters} & \textbf{Mean} & \textbf{Variance} & \textbf{What it models} \\
		\midrule
		\ding{72} Bernoulli 
		& $p$ 
		& $p$ 
		& $p(1-p)$ 
		& Single success/failure trial \\
		
		\ding{72} Binomial 
		& $n,p$ 
		& $np$ 
		& $np(1-p)$ 
		& Number of successes in $n$ trials \\
		
		\ding{72} Geometric$^{*}$ 
		& $p$ 
		& $(1-p)/p$ 
		& $(1-p)/p^2$ 
		& Failures before first success \\
		
		Negative Binomial 
		& $r,p$ 
		& $r(1-p)/p$ 
		& $r(1-p)/p^2$ 
		& Failures before the $r$th success \\
		
		Hypergeometric 
		& $N,K,n$ 
		& $nK/N$ 
		& $n(K/N)(1-K/N)\tfrac{N-n}{N-1}$ 
		& Sampling without replacement \\
		
		\ding{72} Poisson 
		& $\lambda$ 
		& $\lambda$ 
		& $\lambda$ 
		& Count of events in fixed time/space \\
		
		\ding{72} Exponential 
		& $\lambda$ 
		& $1/\lambda$ 
		& $1/\lambda^2$ 
		& Time to next event in Poisson process \\
		
		\ding{72} Continuous Uniform $(a,b)$ 
		& $a<b$ 
		& $(a+b)/2$ 
		& $(b-a)^2/12$ 
		& Constant density on an interval \\
		
		Normal 
		& $\mu,\sigma^2$ 
		& $\mu$ 
		& $\sigma^2$ 
		& Real-valued measurements with symmetric variation (CLT/averages) \\
		\bottomrule
	\end{tabular}
	
	\footnotesize $^{*}$Geometric counts failures before first success; trials-to-success version has mean $1/p$.
	\caption{Reference table of distributions studied in MATH163. Distributions marked with \ding{72} require memorised mean/variance.}
\end{sidewaystable}

%============================================================
\section{Probability Tools}

\textbf{Total probability:}
\[
\Pr(B)=\sum_i \Pr(B\mid A_i)\Pr(A_i).
\]
\textbf{Bayes' rule:}
\[
\Pr(A\mid B)=\frac{\Pr(B\mid A)\Pr(A)}{\Pr(B)}.
\]

%============================================================
\section*{Bias}

An estimator $\hat{\theta}$ of a parameter $\theta$ has \textbf{bias}
\[
\text{Bias}(\hat{\theta}) = \mathbb{E}[\hat{\theta}] - \theta.
\]

\begin{itemize}
	\item \textbf{Unbiased estimator}: $\mathbb{E}[\hat{\theta}] = \theta$.
	\item \textbf{Biased estimator}: expectation differs from the true parameter.
	\item Sample mean $\bar X_n$ is unbiased for $\mu$.
	\item Sample variance $s^2 = \frac{1}{n-1}\sum (x_i - \bar x)^2$ is unbiased for $\sigma^2$.
\end{itemize}

%============================================================
\section*{Consistency}

An estimator $\hat{\theta}_n$ is \textbf{consistent} for $\theta$ if
\[
\hat{\theta}_n \xrightarrow{p} \theta,
\]
i.e.\ it converges in probability to the true parameter as $n \to \infty$.

\begin{itemize}
	\item Equivalent: $\mathbb{E}[\hat{\theta}_n] \to \theta$ and $\mathrm{Var}(\hat{\theta}_n) \to 0$.
	\item $\bar X_n$ is consistent for $\mu$.
	\item Many biased estimators can still be consistent.
\end{itemize}

%============================================================
\section*{Limit laws}

\textbf{Law of Large Numbers (LLN):}
\[
\bar X_n \xrightarrow{p} \mu.
\]
The sample mean stabilises around the population mean as $n$ increases.

\textbf{Central Limit Theorem (CLT):}
\[
\sqrt{n}(\bar X_n - \mu) \xrightarrow{d} N(0, \sigma^2).
\]
For large $n$, $\bar X_n$ is approximately normal even if the data are not.

These justify large-sample $z$-tests and confidence intervals.

%============================================================
\section{Statistical Tests: z and t}

\begin{keybox}
Hypotheses: two‑tail $H_0: \mu=\mu_0$; one‑tail $H_0: \mu\le \mu_0$ or $H_0: \mu\ge \mu_0$.\\
Test statistic:
\[
\frac{\bar X-\mu_0}{\mathrm{SE}}.
\]
\textbf{SE for $z$}: $\sigma/\sqrt{n}$.\\
\textbf{SE for $t$}: $s/\sqrt{n}$ (df $=n-1$).\\
Use $z$ if $\sigma$ known or $n$ large; $t$ otherwise.\\
Small samples ($n<30$) require approximate normality. 
\end{keybox}

Raw data formulas:
\[
\bar x=\frac1n\sum x_i,
\qquad
s^2=\frac1{n-1}\sum (x_i-\bar x)^2.
\]
For frequency data:
\[
\bar x=\frac{1}{n}\sum_j f_j v_j,
\qquad
s^2=\frac1{n-1}\sum_j f_j (v_j-\bar x)^2=\frac n{n-1}\left[\overline{x^2}-(\bar x)^2\right].
\]

%============================================================
\section{Statistical Tests: $\chi^2$}
Statistic:
\[
\chi^2=\sum_i \frac{(O_i-E_i)^2}{E_i}.
\]
Always right‑tail.

Degrees of freedom:
\begin{itemize}
    \item Goodness of fit with $k$ categories: $k-1-$ parameters estimated.
    \item Independence/homogeneity ($r\times c$): $(r-1)(c-1)$ \text{ (no parameters to be estimated).}
\end{itemize}
Conditions: independent samples; expected counts usually $\ge5$.

%============================================================
%\clearpage
\section*{Typical exam format}

\begin{itemize}
	\item Section A contains 10 multiple choice questions, each worth 6 marks. Section A carries 60\% of the available marks. 
	\item In Section A, unless stated otherwise, 3 marks are awarded for the correct option and 3 marks for justification/working.
	\item Section B contains 2 longer questions worth 20 marks each, and carries 40\% of the available marks.
\end{itemize}

\vspace{0.5em}
\textbf{Typical MCQ topics include:}
\begin{itemize}
	\item Identifying a data type
	\item Applying basic probability rules
	\item Counting ordered / unordered combinations
	\item Transforming a discrete random variable
	\item Mean/variance of transformations or combinations of RVs
	\item Covariance; identifying (in)dependence
	\item Identifying the correct RV model in context
	\item Bias of estimators; identifying unbiased estimators
	\item Limit laws; confidence intervals — selecting correct statements
	\item Properties of hypothesis tests, e.g.\ degrees of freedom
\end{itemize}

\vspace{0.5em}
\textbf{QB1: Given a PDF or CDF, tasks may include:}
\begin{itemize}
	\item Determining the normalisation constant
	\item Giving the CDF or PDF (including piecewise form)
	\item Computing the mean $\mathbb E[X]$
	\item Finding the median and commenting on skewness
	\item Computing the variance
	\item Applying a monotone transformation and giving the new PDF
	\item Giving summary statistics of the transformed RV
\end{itemize}

\vspace{0.5em}
\textbf{QB2: Statistical testing. Tasks may include:}
\begin{itemize}
	\item Identifying the correct test ($z$, $t$, or $\chi^2$; one‐tail or two‐tail)
	\item Formulating the null and alternative hypotheses
	\item Determining the degrees of freedom
	\item Computing the test statistic or the $p$–value
	\item Deciding whether to reject $H_0$
	\item Interpreting the decision in context
	\item Identifying assumptions or limitations of the test
\end{itemize}

%============================================================
\end{document}
